\section{Integration}

\subsection{Integration of differential forms}

When we think of integration, usually our theories are built on the integration of real-valued functions.
This is not suitable in the context of manifolds in general.
At least not in a way that would behave nicely.
We would want our integrals to be independent of choice of coordinates, but this is easily seen to be impossible.

Take for example a closed ball $C\subseteq\bR^n$ and the constant function $f\colon C\to\bR^n$ given by $f(x)=1$.
Any sensible theory of integration would mandate that the integral of $f$ is equal to the volume of $C$:
$$ \int_Cf = {\rm Vol}(C) . $$
But this is clearly not coordinate-independent, as $C$ is diffeomorphic to a larger closed ball with a larger volume say.

As we will see, we can solve this via differential forms.

\bdefn

    A {\emphcolor domain of integration} in $\bR^n$ is a bounded subset with measure zero boundary.
    For a domain of integration $D\subseteq\bR^n$, let $\omega$ be a continuous $n$-form on $\cl D$.
    Such an $n$-form can be written $\omega=f\d x^1\wedge\cdots\wedge\d x^n$ for a continuous $f\colon\cl D\to\bR$.

    Then we define $\omega$'s {\emphcolor integral} over $D$ as
    $$ \int_D\omega = \int_Df\d V, $$
    where $\d V$ denotes integration relative to the Lebesgue measure of $\bR^n$.
    We can write this more clearly as
    $$ \int_Df\d x^1\wedge\cdots\wedge\d x^n = \int_Df\d x^1\cdots\d x^n . $$
    So we can view integration as ``removing the wedges''.

    More generally, let $U$ be an open subset of $\bR^n$ or $\bH^n$ and $\omega$ is a compactly supported $n$-form
    on $U$, then define
    $$ \int_U\omega = \int_D\omega $$
    where $D\subseteq\bR^n$ or $\bH^n$ is any domain of integration (e.g.\ a rectangle) containing $\supp\omega$
    where we potentially extend $\omega$ to be zero outside of $U$.

\edefn

This definition does not depend on choice of domain of integration.
This is because for any two domains of integration, $\omega$ is zero outside of their intersection.

\bprop

    Suppose $D,E$ are open domains of integration in $\bR^n$ or $\bH^n$, and $G\colon\cl D\to\cl E$ is a smooth map which
    restricts to an orientation-preserving or orientation-reversing diffeomorphism from $D$ to $E$.
    Then for an $n$-form $\omega$ on $\cl E$, we have
    $$ \int_DG^*\omega = \pm\int_E\omega , $$
    where the sign is flipped when $G$ is orientation-reversing.

\eprop

\bproof

    Denote the standard coordinates of $E$ by $(y^i)$ and those of $D$ by $(x^i)$.
    If $G$ is orientation preserving, let $\omega=f\d y^1\wedge\cdots\wedge\d y^n$, then we have
    $$ \multline{
        \int_E\omega = \int_Ef\d V = \int_D(f\circ G)\abs{\det DG}\d V = \int_D(f\circ G)\det(DG)\d V\cr
        = \int_D(f\circ G)\det(DG)\d x^1\wedge\cdots\wedge\d x^n = \int_DG^*\omega .
    } $$
    If $G$ is orientation-reversing $\abs{\det DG}=-\det(DG)$ and the sign flips.
    \qed

\eproof

\blemm

    Suppose $U$ is an open subset of $\bR^n$ or $\bH^n$, and $K$ is a compact subset of $U$.
    Then there is an open domain of integration $D$ such that $K\subseteq D\subseteq\cl D\subseteq U$.

\elemm

\bproof

    For each $p\in K$ take a open ball or half-open ball whose closure is contained in $U$.
    As $K$ is compact, this has a finite subcovering of balls; let $D$ be their union.
    $D$ is then bound, and its boundary is contained in the union of the boundaries of the balls.
    As a finite union of measure zero sets, the boundary of $D$ has measure zero.
    \qed

\eproof

So we can now generalize our above proposition to compactly supported differential $n$-forms defined on any open subset,
not just a domain of integration.

\bprop

    If $U,V$ are open subsets of $\bR^n$ or $\bH^n$ and $G\colon U\to V$ is an orientation-preserving or reversing
    diffeomorphism, for any compactly supported $n$-form $\omega$ on $V$,
    $$ \int_V\omega = \pm\int_UG^*\omega . $$
    The sign depends on if $G$ preserves or reverses orientation (positive and negative respectively).

\eprop

\bproof

    Take an open domain of integration $E\subseteq V$ such that $\supp\omega\subseteq E$.
    Then $D=G^{-1}E$ is open, and since diffeomorphisms preserve measure-zero-ness and boundaries, we see that
    $D$ is a domain of integration.
    Then we have by the previous proposition,
    $$ \int_V\omega = \int_E\omega = \pm\int_DG^*\omega = \pm\int_UG^*\omega . \qed $$

\eproof

We now extend these definitions to arbitrary oriented smooth manifolds.

Let $M$ be an oriented smooth $n$-manifold, and $\omega$ an $n$-form on $M$.
First suppose that $\omega$ is compactly supported in the domain of a single smooth chart $(U,\phi)$.
Such a chart is either positively or negatively oriented (if connected, otherwise require it to have a consistent orientation),
then we define
$$ \int_M\omega = \pm\int_{\phi(U)}(\phi^{-1})^*\omega, $$
where the right integral is done in Euclidean (half) space, and the sign depends on the orientation of the chart
(positive when oriented, negative when negatively oriented).

\blemm

    The above definition is independent on choice of smooth chart containing $\supp\omega$.

\elemm

\bproof

    If $(U,\phi)$ and $(\tilde U,\tilde\phi)$ both contain $\supp\omega$ then $\tilde\phi\circ\phi^{-1}$
    is a diffeomorphism $\phi(U\cap\tilde U)$ to $\tilde\phi(U\cap\tilde U)$.
    It is orientation-preserving when both are consistently oriented, and negatively oriented otherwise.

    If they are consistently oriented, then
    $$ \multline{
        \int_{\tilde\phi(U)}(\tilde\phi^{-1})^*\omega = \int_{\tilde\phi(U\cap\tilde U)}(\tilde\phi^{-1})^*\omega
        = \int_{\phi(U\cap\tilde U)}(\tilde\phi\circ\phi^{-1})^*\circ(\tilde\phi^{-1})^*\omega\cr
        = \int_{\phi(U\cap\tilde U)}(\phi^{-1})^*\omega = \int_{\phi(U)}(\phi^{-1})^*\omega .
    } $$
    If they have different orientations then the above computation swaps signs, which is what we need by definition.
    \qed

\eproof

Now suppose that $\omega$ is a compactly supported $n$-form on all of $M$.
Let $U_i$ be a finite cover of $\supp\omega$, and let $\psi_i$ be a partition of unity subordinate to it.
Define $\omega$'s integral as
$$ \int_M\omega = \sum_i\int_M\psi_i\omega . $$
Since $\psi_i\omega$ is compactly supported (as both $\omega$ and $\psi_i$ are), this is well-defined, but it is not
immediately clear why this would be independent on choice of partition of unity.

\bnote

    One may ask why we even allow for negatively oriented charts, and instead just use positively oriented charts.
    The reason is that for a one-dimensional manifold with boundary, we cannot ensure the existence of positively
    oriented boundary charts.

\enote

\bprop

    The definition of $\int_M\omega$ does not depend on choice of open cover or partition of unity.

\eprop

\bproof

    Suppose $\tilde U_i$ is another finite open cover of $\supp\omega$ with $\tilde\psi_j$ a subordinate partition
    of unity.
    Then we have for each $i$,
    $$ \int_M\psi_i\omega = \int_M\parens{\sum_j\tilde\psi_j}\psi_i\omega = \sum_j\int_M\tilde\psi_j\psi_i\omega . $$
    Summing over $i$ we see
    $$ \sum_i\int_M\psi_i\omega = \sum_{i,j}\int_M\tilde\psi_j\psi_i\omega . $$
    Each term in the sum is compactly supported in a single smooth chart (e.g.\ $U_i$), so each term is uniquely defined.
    The same process shows
    $$ \sum_j\int_M\tilde\psi_j\omega = \sum_{i,j}\int_M\tilde\psi_j\psi_i\omega , $$
    proving equality.
    \qed

\eproof

For the zero-dimensional case, a compactly supported $0$-form is just a real valued function $f\colon M\to\bR$
which is nonzero on finitely many points.
So we define
$$ \int_Mf = \sum_{p\in M}\pm f(p), $$
where the sign is determined by choice of orientation at $p$.

If $S\subseteq M$ is an immersed oriented $k$-submanifold (with induced orientation), and $\omega$ is a $k$-form
on $M$ whose restriction to $S$ is compactly supported, we define
$$ \int_S\omega = \int_S\iota^*_S\omega, $$
where $\iota^*_S$ is inclusion.
In particular for compact oriented $M$, and an $(n-1)$-form $\omega$ on $M$,
$$ \int_{\partial M}\omega = \int_{\partial M}\iota^*_{\partial M}\omega , $$
where $\partial M$ has the induced (Stokes) orientation.

\bnote

    It is possible to extend integration to non-compactly-supported forms, but this introduces issues of convergence.
    Such definitions are useful, but unnecessary for our purposes.

\enote

\bprop

    Suppose $M,N$ are non-empty oriented smooth $n$-manifolds, $\omega,\eta$ are compactly supported $n$-forms on $M$.
    The integral has the following properties.

    \benum
        \item {\emphcolor Linearity}: for $a,b\in\bR$:
        $$ \int_M a\omega+b\eta = a\int_M\omega + b\int_M\eta . $$
        \item {\emphcolor Orientation reversal}: if $-M$ denotes $M$ with the opposite orientation, then
        $$ \int_{-M}\omega = -\int_M\omega . $$
        \item {\emphcolor Positivity}: if $\omega$ is a positively oriented orientation form, $\int_M\omega>0$.
        \item {\emphcolor Diffeomorphism invariance}: if $F\colon N\to M$ is a orientation-preserving
        or reversing diffeomorphism,
        $$ \int_M\omega = \pm\int_NF^*\omega , $$
        with positive sign when $F$ preserves orientation, and negative sign when it reverses it.
    \eenum

\eprop

\bproof

    (1) is clear at each step in the definition (in Euclidean space, then this implies it for forms supported in a single
    chart, and this implies it in general).
    (2) is an immediate consequence of (4) via $F=\id_M$ as it reverses orientation.

    For (3) if $\omega$ is positively oriented and $(U,\phi)$ is a positively oriented smooth chart, $(\phi^{-1})^*\omega$
    is a positive form on $\phi(U)$, so $f\d x^1\wedge\cdots\wedge\d x^n$ for $f$ positive.
    Then its integral over $U$ is positive.
    For negatively oriented charts, $f$ is negative, but we swap signs so the integral remains positive.
    Then for a partition of unity, each term in the sum is nonnegative, and there must be at least one strictly positive term.

    For (4) it is sufficient to prove this for $\omega$ supported in a single chart, since a partition of unity
    in $M$ maps to a partition of unity of $N$ via $F$.
    So suppose $(U,\phi)$ is positively oriented with $\supp\omega\subseteq U$.
    If $F$ is positively oriented, $(F^{-1}U,\phi\circ F)$ is an oriented smooth atlas for $N$ whose domain
    contains $\supp F^*\omega$.
    The result then follows from diffeomorphism invariance of the Euclidean integral.
    The other cases are handled similarly.
    \qed

\eproof

\subsection{Stoke's theorem}

\bthrm

    Let $M$ be an oriented $n$-manifold with boundary, and $\omega$ a compactly supported smooth $(n-1)$-form on $M$.
    Then
    $$ \int_M\d\omega = \int_{\partial M}\omega, $$
    where $\partial M$ is understood to have the Stokes orientation.

\ethrm

\bproof

    First we consider the case that $M=\bH^n$.
    Since $\omega$ is compactly supported, its support is contained within a rectangle $A=[-R,R]^{n-1}\times[0,R]$.
    Writing $\omega$ in terms of the standard coordinates,
    $$ \omega = \sum_{i=1}^n\omega_i\d x^1\wedge\cdots\wedge\hat{\d x^i}\wedge\cdots\wedge\d x^n . $$
    Thus,
    $$ \d\omega = \sum_{i=1}^n\d\omega_i\wedge\d x^1\wedge\cdots\wedge\hat{\d x^i}\wedge\cdots\wedge\d x^n =
    \sum_{i,j}\pdvof{\omega_i}{x^j}\d x^j\wedge\d x^1\wedge\cdots\wedge\hat{\d x^i}\wedge\cdots\wedge\d x^n =
    \sum_i(-1)^{i-1}\pdvof{\omega_i}{x^i}\d x^1\wedge\cdots\wedge\d x^i\wedge\cdots\wedge\d x^n . $$
    Then we compute,
    $$ \int_{\bH^n}\d\omega =
    \sum_{i=1}^n(-1)^{i-1}\int_A\pdvof{\omega_i}{x^i}\d x^1\wedge\cdots\wedge\d x^i\wedge\cdots\wedge\d x^n =
    \sum_{i=1}^n(-1)^{i-1}\int_0^R\int_0^R\cdots\int_0^R\pdvof{\omega_i}{x^i}\d x^1\cdots\d x^n . $$
    By Fubini we can change the order of integration as to integrate the $i$th term first.
    When $i\neq n$

\eproof

